%%%%%%%%%%%%%%%%%%%%%%%%%%%%%%%%%%%%%%%%%%%%%%%%%%%%%%%%%%%%
% 7.10 CENTRAL LIMIT THEOREMS
% The Composition of Advanced Mathematics
%%%%%%%%%%%%%%%%%%%%%%%%%%%%%%%%%%%%%%%%%%%%%%%%%%%%%%%%%%%%

\section{Central Limit Theorems}
\label{sec:central-limit-theorems}

\prerequisite{
Random variables, expectation, variance, independence, identical
distributions, laws of large numbers, probability transforms,
characteristic functions, standardization, normal distributions, and
basic asymptotic notation.
}

\learningobjectives{
By the end of this section, the reader should be able to:
\begin{itemize}
    \item explain the purpose of a central limit theorem;
    \item distinguish law-of-large-numbers behavior from
          central-limit behavior;
    \item standardize sums and sample means;
    \item state the classical i.i.d. central limit theorem;
    \item apply the CLT to sums of independent random variables;
    \item apply the CLT to sample means;
    \item understand convergence in distribution;
    \item compute approximate probabilities using a normal limit;
    \item approximate binomial distributions by normal distributions;
    \item use continuity corrections;
    \item understand the de Moivre--Laplace theorem;
    \item recognize the role of finite variance;
    \item understand characteristic-function proofs at an introductory
          level;
    \item recognize Lindeberg--Lévy, Lyapunov, and Lindeberg--Feller
          forms of the CLT;
    \item understand triangular arrays at an introductory level;
    \item interpret the Berry--Esseen theorem as a convergence-rate
          result;
    \item distinguish finite-sample normality from asymptotic normality;
    \item recognize when normal approximations may be poor;
    \item understand standardized sums of nonnormal variables;
    \item recognize multivariate central limit theorems;
    \item connect the CLT with statistics, Monte Carlo methods,
          confidence intervals, hypothesis testing, and stochastic
          processes.
\end{itemize}
}

%%%%%%%%%%%%%%%%%%%%%%%%%%%%%%%%%%%%%%%%%%%%%%%%%%%%%%%%%%%%
\subsection{Why Central Limit Theorems Matter}
\label{subsec:why-central-limit-theorems}
%%%%%%%%%%%%%%%%%%%%%%%%%%%%%%%%%%%%%%%%%%%%%%%%%%%%%%%%%%%%

The law of large numbers explains why averages stabilize.

It tells us that

\[
\overline{X}_n
\]

approaches

\[
\mu.
\]

However, it does not describe the shape of the random error

\[
\overline{X}_n-\mu.
\]

Central limit theorems answer this second question.

They describe the distribution of suitably standardized sums and
averages.

The central principle is

\[
\boxed{
\text{many small independent contributions}
\longrightarrow
\text{approximately normal fluctuations}.
}
\]

%%%%%%%%%%%%%%%%%%%%%%%%%%%%%%%%%%%%%%%%%%%%%%%%%%%%%%%%%%%%
\subsection{Sample Sums and Sample Means}
\label{subsec:clt-sample-sums-means}
%%%%%%%%%%%%%%%%%%%%%%%%%%%%%%%%%%%%%%%%%%%%%%%%%%%%%%%%%%%%

Let

\[
X_1,X_2,\ldots
\]

be random variables.

Define

\[
\boxed{
S_n
=
X_1+\cdots+X_n.
}
\]

The sample mean is

\[
\boxed{
\overline{X}_n
=
\frac{S_n}{n}.
}
\]

Suppose the variables are i.i.d. with

\[
\mathbb{E}[X_i]=\mu
\]

and

\[
\operatorname{Var}(X_i)=\sigma^2.
\]

Then

\[
\boxed{
\mathbb{E}[S_n]
=
n\mu
}
\]

and

\[
\boxed{
\operatorname{Var}(S_n)
=
n\sigma^2.
}
\]

%%%%%%%%%%%%%%%%%%%%%%%%%%%%%%%%%%%%%%%%%%%%%%%%%%%%%%%%%%%%
\subsection{Standardizing a Sum}
\label{subsec:standardizing-sum-clt}
%%%%%%%%%%%%%%%%%%%%%%%%%%%%%%%%%%%%%%%%%%%%%%%%%%%%%%%%%%%%

The standard deviation of \(S_n\) is

\[
\sqrt{n}\sigma.
\]

Therefore the standardized sum is

\[
\boxed{
Z_n
=
\frac{
S_n-n\mu
}{
\sigma\sqrt{n}
}.
}
\]

This random variable has mean

\[
0
\]

and variance

\[
1.
\]

The central limit theorem describes the limiting distribution of
\(Z_n\).

%%%%%%%%%%%%%%%%%%%%%%%%%%%%%%%%%%%%%%%%%%%%%%%%%%%%%%%%%%%%
\subsection{Standardizing the Sample Mean}
\label{subsec:standardizing-sample-mean-clt}
%%%%%%%%%%%%%%%%%%%%%%%%%%%%%%%%%%%%%%%%%%%%%%%%%%%%%%%%%%%%

Since

\[
S_n=n\overline{X}_n,
\]

we have

\[
\frac{
S_n-n\mu
}{
\sigma\sqrt{n}
}
=
\frac{
\sqrt{n}
(\overline{X}_n-\mu)
}{
\sigma
}.
\]

Thus the standardized sample mean is

\[
\boxed{
Z_n
=
\frac{
\overline{X}_n-\mu
}{
\sigma/\sqrt{n}
}.
}
\]

The quantity

\[
\boxed{
\frac{\sigma}{\sqrt{n}}
}
\]

is the standard deviation of the sample mean and is commonly called its
standard error.

%%%%%%%%%%%%%%%%%%%%%%%%%%%%%%%%%%%%%%%%%%%%%%%%%%%%%%%%%%%%
\subsection{Convergence in Distribution}
\label{subsec:convergence-in-distribution-clt}
%%%%%%%%%%%%%%%%%%%%%%%%%%%%%%%%%%%%%%%%%%%%%%%%%%%%%%%%%%%%

\begin{definitionbox}{Convergence in Distribution}
A sequence of random variables \(X_n\) converges in distribution to a
random variable \(X\) if

\[
\boxed{
F_{X_n}(x)
\to
F_X(x)
}
\]

at every point \(x\) where \(F_X\) is continuous.
\end{definitionbox}

We write

\[
\boxed{
X_n
\xrightarrow{d}
X.
}
\]

The symbol

\[
\xrightarrow{d}
\]

is read

\[
\text{``converges in distribution.''}
\]

%%%%%%%%%%%%%%%%%%%%%%%%%%%%%%%%%%%%%%%%%%%%%%%%%%%%%%%%%%%%
\subsection{Meaning of Convergence in Distribution}
\label{subsec:meaning-convergence-distribution}
%%%%%%%%%%%%%%%%%%%%%%%%%%%%%%%%%%%%%%%%%%%%%%%%%%%%%%%%%%%%

Convergence in distribution means that the distribution functions of
\(X_n\) approach the distribution function of \(X\).

It is weaker than convergence in probability.

Recall the implication

\[
\boxed{
X_n
\xrightarrow{P}
X
\Longrightarrow
X_n
\xrightarrow{d}
X.
}
\]

The converse is not generally true.

%%%%%%%%%%%%%%%%%%%%%%%%%%%%%%%%%%%%%%%%%%%%%%%%%%%%%%%%%%%%
\subsection{Classical Central Limit Theorem}
\label{subsec:classical-central-limit-theorem}
%%%%%%%%%%%%%%%%%%%%%%%%%%%%%%%%%%%%%%%%%%%%%%%%%%%%%%%%%%%%

\begin{theorem}[Lindeberg--Lévy Central Limit Theorem]
Let

\[
X_1,X_2,\ldots
\]

be independent and identically distributed random variables with

\[
\mathbb{E}[X_i]=\mu
\]

and

\[
0<
\operatorname{Var}(X_i)
=
\sigma^2
<
\infty.
\]

Then

\[
\boxed{
\frac{
S_n-n\mu
}{
\sigma\sqrt{n}
}
\xrightarrow{d}
N(0,1).
}
\]

Equivalently,

\[
\boxed{
\frac{
\sqrt{n}
(\overline{X}_n-\mu)
}{
\sigma
}
\xrightarrow{d}
N(0,1).
}
\]
\end{theorem}

%%%%%%%%%%%%%%%%%%%%%%%%%%%%%%%%%%%%%%%%%%%%%%%%%%%%%%%%%%%%
\subsection{CDF Form of the CLT}
\label{subsec:cdf-form-clt}
%%%%%%%%%%%%%%%%%%%%%%%%%%%%%%%%%%%%%%%%%%%%%%%%%%%%%%%%%%%%

The theorem means that for every real \(z\),

\[
\boxed{
P
\left(
\frac{
S_n-n\mu
}{
\sigma\sqrt{n}
}
\leq z
\right)
\to
\Phi(z),
}
\]

where

\[
\Phi(z)
=
P(Z\leq z)
\]

is the standard normal CDF.

Thus large standardized sums behave approximately like a standard
normal variable.

%%%%%%%%%%%%%%%%%%%%%%%%%%%%%%%%%%%%%%%%%%%%%%%%%%%%%%%%%%%%
\subsection{Approximate Distribution of a Sample Sum}
\label{subsec:approx-distribution-sum-clt}
%%%%%%%%%%%%%%%%%%%%%%%%%%%%%%%%%%%%%%%%%%%%%%%%%%%%%%%%%%%%

For sufficiently large \(n\),

\[
\boxed{
S_n
\approx
N
\left(
n\mu,
n\sigma^2
\right).
}
\]

This approximation is not an identity unless the original variables are
normal.

It is an asymptotic approximation.

%%%%%%%%%%%%%%%%%%%%%%%%%%%%%%%%%%%%%%%%%%%%%%%%%%%%%%%%%%%%
\subsection{Approximate Distribution of a Sample Mean}
\label{subsec:approx-distribution-mean-clt}
%%%%%%%%%%%%%%%%%%%%%%%%%%%%%%%%%%%%%%%%%%%%%%%%%%%%%%%%%%%%

For sufficiently large \(n\),

\[
\boxed{
\overline{X}_n
\approx
N
\left(
\mu,
\frac{\sigma^2}{n}
\right).
}
\]

Equivalently,

\[
\boxed{
\overline{X}_n
\approx
N
\left(
\mu,
\left(
\frac{\sigma}{\sqrt{n}}
\right)^2
\right).
}
\]

This approximation is fundamental throughout statistics.

%%%%%%%%%%%%%%%%%%%%%%%%%%%%%%%%%%%%%%%%%%%%%%%%%%%%%%%%%%%%
\subsection{Worked Example: Sample Mean Approximation}
\label{subsec:worked-clt-sample-mean}
%%%%%%%%%%%%%%%%%%%%%%%%%%%%%%%%%%%%%%%%%%%%%%%%%%%%%%%%%%%%

\begin{workedexamplebox}{Approximating a Sample Mean}

Suppose a population has

\[
\mu=50
\]

and

\[
\sigma=12.
\]

Take an independent sample of size

\[
n=36.
\]

Then

\[
\mathbb{E}[\overline{X}]
=
50,
\]

and

\[
\operatorname{SD}(\overline{X})
=
\frac{12}{\sqrt{36}}
=
2.
\]

For sufficiently appropriate conditions,

\[
\begin{answerbox}
\[
\boxed{
\overline{X}
\approx
N(50,2^2).
}
\]
\end{answerbox}

\end{workedexamplebox}

%%%%%%%%%%%%%%%%%%%%%%%%%%%%%%%%%%%%%%%%%%%%%%%%%%%%%%%%%%%%
\subsection{Worked Example: Probability for a Sample Mean}
\label{subsec:worked-clt-sample-mean-probability}
%%%%%%%%%%%%%%%%%%%%%%%%%%%%%%%%%%%%%%%%%%%%%%%%%%%%%%%%%%%%

\begin{workedexamplebox}{Probability Using the CLT}

Continue with

\[
\mu=50,
\qquad
\sigma=12,
\qquad
n=36.
\]

Approximate

\[
P(\overline{X}\leq54).
\]

Standardize:

\[
z
=
\frac{
54-50
}{
12/\sqrt{36}
}.
\]

Thus,

\[
z=2.
\]

Therefore,

\begin{answerbox}
\[
\boxed{
P(\overline{X}\leq54)
\approx
\Phi(2).
}
\]
\end{answerbox}

\end{workedexamplebox}

%%%%%%%%%%%%%%%%%%%%%%%%%%%%%%%%%%%%%%%%%%%%%%%%%%%%%%%%%%%%
\subsection{Exact Normality Versus CLT Normality}
\label{subsec:exact-versus-asymptotic-normality}
%%%%%%%%%%%%%%%%%%%%%%%%%%%%%%%%%%%%%%%%%%%%%%%%%%%%%%%%%%%%

If

\[
X_1,\ldots,X_n
\]

are themselves normal, then

\[
\overline{X}_n
\]

is exactly normal for every \(n\).

If the \(X_i\) are not normal, the CLT says that the standardized mean
becomes approximately normal as \(n\) becomes large.

Thus,

\[
\boxed{
\text{normal population}
\Longrightarrow
\text{exact normal sample mean},
}
\]

whereas

\[
\boxed{
\text{general finite-variance population}
\Longrightarrow
\text{asymptotic normality}.
}
\]

%%%%%%%%%%%%%%%%%%%%%%%%%%%%%%%%%%%%%%%%%%%%%%%%%%%%%%%%%%%%
\subsection{The CLT Does Not Say the Data Become Normal}
\label{subsec:clt-data-not-normal}
%%%%%%%%%%%%%%%%%%%%%%%%%%%%%%%%%%%%%%%%%%%%%%%%%%%%%%%%%%%%

The CLT does not claim that the individual observations

\[
X_i
\]

become normally distributed.

It concerns sums and averages.

A highly skewed population may remain highly skewed.

What becomes approximately normal is

\[
\boxed{
\frac{
S_n-n\mu
}{
\sigma\sqrt{n}
}
}
\]

or equivalently the standardized sample mean.

%%%%%%%%%%%%%%%%%%%%%%%%%%%%%%%%%%%%%%%%%%%%%%%%%%%%%%%%%%%%
\subsection{LLN Versus CLT}
\label{subsec:clt-versus-lln}
%%%%%%%%%%%%%%%%%%%%%%%%%%%%%%%%%%%%%%%%%%%%%%%%%%%%%%%%%%%%

The law of large numbers gives

\[
\boxed{
\overline{X}_n
\to
\mu.
}
\]

The central limit theorem studies

\[
\boxed{
\sqrt{n}
(\overline{X}_n-\mu).
}
\]

The LLN says the error shrinks.

The CLT describes the distributional shape of that shrinking error.

Thus:

\[
\boxed{
\text{LLN}
=
\text{location of the average},
}
\]

while

\[
\boxed{
\text{CLT}
=
\text{shape of the fluctuations}.
}
\]

%%%%%%%%%%%%%%%%%%%%%%%%%%%%%%%%%%%%%%%%%%%%%%%%%%%%%%%%%%%%
\subsection{Why Standardization Is Necessary}
\label{subsec:why-standardization-clt}
%%%%%%%%%%%%%%%%%%%%%%%%%%%%%%%%%%%%%%%%%%%%%%%%%%%%%%%%%%%%

Without centering,

\[
S_n
\]

typically grows like

\[
n\mu.
\]

Without scaling,

\[
S_n-n\mu
\]

typically has standard deviation

\[
\sigma\sqrt{n}.
\]

Therefore the normalization

\[
\boxed{
\frac{
S_n-n\mu
}{
\sigma\sqrt{n}
}
}
\]

removes both deterministic growth and increasing spread.

The result remains on a stable probability scale.

%%%%%%%%%%%%%%%%%%%%%%%%%%%%%%%%%%%%%%%%%%%%%%%%%%%%%%%%%%%%
\subsection{Characteristic-Function Idea of the CLT}
\label{subsec:cf-proof-clt-idea}
%%%%%%%%%%%%%%%%%%%%%%%%%%%%%%%%%%%%%%%%%%%%%%%%%%%%%%%%%%%%

Suppose

\[
X_i
\]

are i.i.d. with mean

\[
0
\]

and variance

\[
1.
\]

Near \(t=0\),

\[
\varphi_X(t)
=
1-\frac{t^2}{2}+o(t^2).
\]

For

\[
Z_n
=
\frac{
X_1+\cdots+X_n
}{
\sqrt{n}
},
\]

independence gives

\[
\varphi_{Z_n}(t)
=
\left[
\varphi_X
\left(
\frac{t}{\sqrt{n}}
\right)
\right]^n.
\]

%%%%%%%%%%%%%%%%%%%%%%%%%%%%%%%%%%%%%%%%%%%%%%%%%%%%%%%%%%%%
\subsection{Characteristic-Function Limit}
\label{subsec:cf-clt-limit}
%%%%%%%%%%%%%%%%%%%%%%%%%%%%%%%%%%%%%%%%%%%%%%%%%%%%%%%%%%%%

Using the expansion,

\[
\varphi_X
\left(
\frac{t}{\sqrt{n}}
\right)
=
1
-
\frac{t^2}{2n}
+
o\left(
\frac1n
\right).
\]

Therefore,

\[
\varphi_{Z_n}(t)
=
\left[
1
-
\frac{t^2}{2n}
+
o\left(
\frac1n
\right)
\right]^n.
\]

As

\[
n\to\infty,
\]

this approaches

\[
\boxed{
e^{-t^2/2}.
}
\]

But

\[
e^{-t^2/2}
\]

is the characteristic function of

\[
N(0,1).
\]

Hence the standardized sum converges in distribution to the standard
normal law.

%%%%%%%%%%%%%%%%%%%%%%%%%%%%%%%%%%%%%%%%%%%%%%%%%%%%%%%%%%%%
\subsection{Role of Lévy's Continuity Theorem}
\label{subsec:levy-role-clt}
%%%%%%%%%%%%%%%%%%%%%%%%%%%%%%%%%%%%%%%%%%%%%%%%%%%%%%%%%%%%

The characteristic-function calculation shows

\[
\varphi_{Z_n}(t)
\to
e^{-t^2/2}.
\]

Lévy's continuity theorem then converts transform convergence into
distributional convergence:

\[
\boxed{
Z_n
\xrightarrow{d}
N(0,1).
}
\]

Thus characteristic functions provide a powerful analytic proof of the
CLT.

%%%%%%%%%%%%%%%%%%%%%%%%%%%%%%%%%%%%%%%%%%%%%%%%%%%%%%%%%%%%
\subsection{de Moivre--Laplace Theorem}
\label{subsec:de-moivre-laplace}
%%%%%%%%%%%%%%%%%%%%%%%%%%%%%%%%%%%%%%%%%%%%%%%%%%%%%%%%%%%%

The earliest central limit theorem arose from the binomial
distribution.

\begin{theorem}[de Moivre--Laplace Theorem]
Let

\[
X_n
\sim
\operatorname{Bin}(n,p)
\]

with

\[
0<p<1.
\]

Then

\[
\boxed{
\frac{
X_n-np
}{
\sqrt{np(1-p)}
}
\xrightarrow{d}
N(0,1).
}
\]
\end{theorem}

Thus the binomial distribution becomes approximately normal after
standardization.

%%%%%%%%%%%%%%%%%%%%%%%%%%%%%%%%%%%%%%%%%%%%%%%%%%%%%%%%%%%%
\subsection{Normal Approximation to the Binomial}
\label{subsec:normal-approx-binomial}
%%%%%%%%%%%%%%%%%%%%%%%%%%%%%%%%%%%%%%%%%%%%%%%%%%%%%%%%%%%%

If

\[
X\sim\operatorname{Bin}(n,p),
\]

then for sufficiently large \(n\),

\[
\boxed{
X
\approx
N
\left(
np,
np(1-p)
\right).
}
\]

The approximate standard deviation is

\[
\boxed{
\sqrt{
np(1-p)
}.
}
\]

%%%%%%%%%%%%%%%%%%%%%%%%%%%%%%%%%%%%%%%%%%%%%%%%%%%%%%%%%%%%
\subsection{When the Binomial Normal Approximation Is Reasonable}
\label{subsec:binomial-normal-conditions}
%%%%%%%%%%%%%%%%%%%%%%%%%%%%%%%%%%%%%%%%%%%%%%%%%%%%%%%%%%%%

A common practical guideline is that both

\[
np
\]

and

\[
n(1-p)
\]

should be reasonably large.

This ensures that the binomial distribution is not too concentrated
against either endpoint.

Exact thresholds vary by context.

The approximation generally improves as \(n\) grows and the
distribution becomes less severely skewed.

%%%%%%%%%%%%%%%%%%%%%%%%%%%%%%%%%%%%%%%%%%%%%%%%%%%%%%%%%%%%
\subsection{Continuity Correction}
\label{subsec:continuity-correction-clt}
%%%%%%%%%%%%%%%%%%%%%%%%%%%%%%%%%%%%%%%%%%%%%%%%%%%%%%%%%%%%

The binomial distribution is discrete, while the normal approximation
is continuous.

A continuity correction adjusts the normal interval by \(0.5\).

For example,

\[
P(X\leq k)
\]

is approximated by

\[
\boxed{
P(Y\leq k+0.5),
}
\]

where \(Y\) is the approximating normal variable.

Similarly,

\[
P(X\geq k)
\]

is approximated by

\[
\boxed{
P(Y\geq k-0.5).
}
\]

%%%%%%%%%%%%%%%%%%%%%%%%%%%%%%%%%%%%%%%%%%%%%%%%%%%%%%%%%%%%
\subsection{Continuity-Correction Table}
\label{subsec:continuity-correction-table}
%%%%%%%%%%%%%%%%%%%%%%%%%%%%%%%%%%%%%%%%%%%%%%%%%%%%%%%%%%%%

For an integer-valued \(X\) approximated by continuous \(Y\),

\[
\begin{array}{c|c}
\text{Discrete Event}
&
\text{Continuous Approximation}
\\
\hline
X=k
&
k-0.5<Y<k+0.5
\\[1mm]
X\leq k
&
Y<k+0.5
\\[1mm]
X<k
&
Y<k-0.5
\\[1mm]
X\geq k
&
Y>k-0.5
\\[1mm]
X>k
&
Y>k+0.5
\\[1mm]
a\leq X\leq b
&
a-0.5<Y<b+0.5
\end{array}
\]

%%%%%%%%%%%%%%%%%%%%%%%%%%%%%%%%%%%%%%%%%%%%%%%%%%%%%%%%%%%%
\subsection{Worked Example: Binomial Normal Approximation}
\label{subsec:worked-binomial-normal-approx}
%%%%%%%%%%%%%%%%%%%%%%%%%%%%%%%%%%%%%%%%%%%%%%%%%%%%%%%%%%%%

\begin{workedexamplebox}{Approximating a Binomial Probability}

Suppose

\[
X
\sim
\operatorname{Bin}(100,0.4).
\]

Then

\[
\mu=np=40
\]

and

\[
\sigma
=
\sqrt{
np(1-p)
}
=
\sqrt{24}.
\]

Approximate

\[
P(X\leq45).
\]

Using the continuity correction,

\[
P(X\leq45)
\approx
P(Y\leq45.5),
\]

where

\[
Y\sim N(40,24).
\]

Standardize:

\[
z
=
\frac{
45.5-40
}{
\sqrt{24}
}.
\]

Therefore,

\begin{answerbox}
\[
\boxed{
P(X\leq45)
\approx
\Phi
\left(
\frac{5.5}{\sqrt{24}}
\right).
}
\]
\end{answerbox}

\end{workedexamplebox}

%%%%%%%%%%%%%%%%%%%%%%%%%%%%%%%%%%%%%%%%%%%%%%%%%%%%%%%%%%%%
\subsection{Worked Example: Exact Integer Probability}
\label{subsec:worked-continuity-single-value}
%%%%%%%%%%%%%%%%%%%%%%%%%%%%%%%%%%%%%%%%%%%%%%%%%%%%%%%%%%%%

\begin{workedexamplebox}{Approximating \(P(X=50)\)}

Suppose

\[
X
\sim
\operatorname{Bin}(100,0.5).
\]

Then

\[
\mu=50
\]

and

\[
\sigma=5.
\]

Approximate

\[
P(X=50).
\]

With continuity correction,

\[
P(X=50)
\approx
P(49.5<Y<50.5)
\]

where

\[
Y\sim N(50,25).
\]

Standardizing gives

\[
-\frac{0.5}{5}
<
Z
<
\frac{0.5}{5}.
\]

Hence,

\begin{answerbox}
\[
\boxed{
P(X=50)
\approx
\Phi(0.1)-\Phi(-0.1).
}
\]
\end{answerbox}

\end{workedexamplebox}

%%%%%%%%%%%%%%%%%%%%%%%%%%%%%%%%%%%%%%%%%%%%%%%%%%%%%%%%%%%%
\subsection{Poisson Normal Approximation}
\label{subsec:poisson-normal-approximation}
%%%%%%%%%%%%%%%%%%%%%%%%%%%%%%%%%%%%%%%%%%%%%%%%%%%%%%%%%%%%

For

\[
X\sim\operatorname{Poisson}(\lambda),
\]

we have

\[
\mathbb{E}[X]=\lambda
\]

and

\[
\operatorname{Var}(X)=\lambda.
\]

For large \(\lambda\),

\[
\boxed{
X
\approx
N(\lambda,\lambda).
}
\]

Because the Poisson variable is discrete, a continuity correction can
again improve the approximation.

%%%%%%%%%%%%%%%%%%%%%%%%%%%%%%%%%%%%%%%%%%%%%%%%%%%%%%%%%%%%
\subsection{Worked Example: Poisson Normal Approximation}
\label{subsec:worked-poisson-normal}
%%%%%%%%%%%%%%%%%%%%%%%%%%%%%%%%%%%%%%%%%%%%%%%%%%%%%%%%%%%%

\begin{workedexamplebox}{Large Poisson Count}

Suppose

\[
X\sim\operatorname{Poisson}(64).
\]

Approximate

\[
P(X\leq70).
\]

Use

\[
Y\sim N(64,64).
\]

The standard deviation is

\[
8.
\]

Applying continuity correction,

\[
P(X\leq70)
\approx
P(Y\leq70.5).
\]

Therefore,

\[
z
=
\frac{
70.5-64
}{8}.
\]

\begin{answerbox}
\[
\boxed{
P(X\leq70)
\approx
\Phi
\left(
\frac{6.5}{8}
\right).
}
\]
\end{answerbox}

\end{workedexamplebox}

%%%%%%%%%%%%%%%%%%%%%%%%%%%%%%%%%%%%%%%%%%%%%%%%%%%%%%%%%%%%
\subsection{CLT for Sample Proportions}
\label{subsec:clt-sample-proportions}
%%%%%%%%%%%%%%%%%%%%%%%%%%%%%%%%%%%%%%%%%%%%%%%%%%%%%%%%%%%%

Let

\[
X_1,\ldots,X_n
\]

be independent Bernoulli\((p)\) variables.

Define the sample proportion

\[
\boxed{
\widehat{p}
=
\frac1n
\sum_{i=1}^{n}X_i.
}
\]

Then

\[
\mathbb{E}[\widehat{p}]
=
p
\]

and

\[
\operatorname{Var}(\widehat{p})
=
\frac{
p(1-p)
}{n}.
\]

The CLT gives

\[
\boxed{
\frac{
\widehat{p}-p
}{
\sqrt{
p(1-p)/n
}
}
\xrightarrow{d}
N(0,1).
}
\]

%%%%%%%%%%%%%%%%%%%%%%%%%%%%%%%%%%%%%%%%%%%%%%%%%%%%%%%%%%%%
\subsection{Approximate Distribution of a Sample Proportion}
\label{subsec:sample-proportion-approx}
%%%%%%%%%%%%%%%%%%%%%%%%%%%%%%%%%%%%%%%%%%%%%%%%%%%%%%%%%%%%

For large \(n\),

\[
\boxed{
\widehat{p}
\approx
N
\left(
p,
\frac{
p(1-p)
}{n}
\right).
}
\]

This approximation becomes fundamental in confidence intervals and
hypothesis tests for proportions.

%%%%%%%%%%%%%%%%%%%%%%%%%%%%%%%%%%%%%%%%%%%%%%%%%%%%%%%%%%%%
\subsection{Worked Example: Sample Proportion}
\label{subsec:worked-sample-proportion-clt}
%%%%%%%%%%%%%%%%%%%%%%%%%%%%%%%%%%%%%%%%%%%%%%%%%%%%%%%%%%%%

\begin{workedexamplebox}{Approximating a Proportion}

Suppose

\[
p=0.6
\]

and

\[
n=400.
\]

Then

\[
\mathbb{E}[\widehat{p}]
=
0.6.
\]

Also,

\[
\operatorname{SD}(\widehat{p})
=
\sqrt{
\frac{
(0.6)(0.4)
}{400}
}.
\]

Therefore,

\begin{answerbox}
\[
\boxed{
\widehat{p}
\approx
N
\left(
0.6,
\frac{0.24}{400}
\right).
}
\]
\end{answerbox}

\end{workedexamplebox}

%%%%%%%%%%%%%%%%%%%%%%%%%%%%%%%%%%%%%%%%%%%%%%%%%%%%%%%%%%%%
\subsection{How Large Must \(n\) Be?}
\label{subsec:how-large-n-clt}
%%%%%%%%%%%%%%%%%%%%%%%%%%%%%%%%%%%%%%%%%%%%%%%%%%%%%%%%%%%%

There is no universal sample size that guarantees a good normal
approximation for every distribution.

The quality of the approximation depends on features such as:

\[
\text{skewness},
\]

\[
\text{tail heaviness},
\]

\[
\text{presence of outliers},
\]

and

\[
\text{existence of moments}.
\]

A nearly symmetric population may require only a modest sample size.

A highly skewed or heavy-tailed population may require much larger
samples.

%%%%%%%%%%%%%%%%%%%%%%%%%%%%%%%%%%%%%%%%%%%%%%%%%%%%%%%%%%%%
\subsection{When the CLT Can Be Slow}
\label{subsec:slow-clt-convergence}
%%%%%%%%%%%%%%%%%%%%%%%%%%%%%%%%%%%%%%%%%%%%%%%%%%%%%%%%%%%%

If the parent distribution is extremely skewed, the distribution of the
sample mean may remain noticeably skewed for moderate \(n\).

If the distribution has very heavy tails, rare extreme values may
dominate the sample sum.

Thus the phrase

\[
\boxed{
\text{``\(n\) is large''}
}
\]

must be interpreted relative to the underlying distribution.

%%%%%%%%%%%%%%%%%%%%%%%%%%%%%%%%%%%%%%%%%%%%%%%%%%%%%%%%%%%%
\subsection{Finite Variance Is Essential in the Classical CLT}
\label{subsec:finite-variance-clt}
%%%%%%%%%%%%%%%%%%%%%%%%%%%%%%%%%%%%%%%%%%%%%%%%%%%%%%%%%%%%

The classical i.i.d. CLT requires

\[
0<\sigma^2<\infty.
\]

If the variance is infinite, the normalization

\[
\sqrt{n}
\]

may no longer be appropriate.

Other limiting distributions, including stable laws, may arise.

Thus,

\[
\boxed{
\text{finite variance is a structural assumption of the classical CLT}.
}
\]

%%%%%%%%%%%%%%%%%%%%%%%%%%%%%%%%%%%%%%%%%%%%%%%%%%%%%%%%%%%%
\subsection{Cauchy Distribution as a Counterexample}
\label{subsec:cauchy-clt-counterexample}
%%%%%%%%%%%%%%%%%%%%%%%%%%%%%%%%%%%%%%%%%%%%%%%%%%%%%%%%%%%%

Let

\[
X_1,\ldots,X_n
\]

be independent standard Cauchy random variables.

Their variance does not exist.

Indeed, the average

\[
\frac{
X_1+\cdots+X_n
}{n}
\]

has the same Cauchy distribution as each \(X_i\).

Therefore the standard normal CLT does not apply.

This shows why moment assumptions cannot simply be ignored.

%%%%%%%%%%%%%%%%%%%%%%%%%%%%%%%%%%%%%%%%%%%%%%%%%%%%%%%%%%%%
\subsection{Lyapunov Central Limit Theorem}
\label{subsec:lyapunov-clt}
%%%%%%%%%%%%%%%%%%%%%%%%%%%%%%%%%%%%%%%%%%%%%%%%%%%%%%%%%%%%

The variables in a CLT need not always be identically distributed.

Suppose

\[
X_1,X_2,\ldots
\]

are independent with

\[
\mu_k
=
\mathbb{E}[X_k]
\]

and variances

\[
\sigma_k^2.
\]

Define

\[
s_n^2
=
\sum_{k=1}^{n}\sigma_k^2.
\]

%%%%%%%%%%%%%%%%%%%%%%%%%%%%%%%%%%%%%%%%%%%%%%%%%%%%%%%%%%%%
\subsection{Lyapunov Condition}
\label{subsec:lyapunov-condition}
%%%%%%%%%%%%%%%%%%%%%%%%%%%%%%%%%%%%%%%%%%%%%%%%%%%%%%%%%%%%

A standard Lyapunov condition requires that for some

\[
\delta>0,
\]

\[
\boxed{
\frac{
1
}{
s_n^{2+\delta}
}
\sum_{k=1}^{n}
\mathbb{E}
\left[
|X_k-\mu_k|^{2+\delta}
\right]
\to0.
}
\]

The Greek letter

\[
\delta
\]

is called ``delta.''

This condition prevents any small collection of variables from
dominating the total variance.

%%%%%%%%%%%%%%%%%%%%%%%%%%%%%%%%%%%%%%%%%%%%%%%%%%%%%%%%%%%%
\subsection{Lyapunov CLT Statement}
\label{subsec:lyapunov-clt-statement}
%%%%%%%%%%%%%%%%%%%%%%%%%%%%%%%%%%%%%%%%%%%%%%%%%%%%%%%%%%%%

Under independence and the Lyapunov condition,

\[
\boxed{
\frac{
\sum_{k=1}^{n}
(X_k-\mu_k)
}{
s_n
}
\xrightarrow{d}
N(0,1).
}
\]

Thus identical distribution is not necessary.

%%%%%%%%%%%%%%%%%%%%%%%%%%%%%%%%%%%%%%%%%%%%%%%%%%%%%%%%%%%%
\subsection{Lindeberg Condition}
\label{subsec:lindeberg-condition}
%%%%%%%%%%%%%%%%%%%%%%%%%%%%%%%%%%%%%%%%%%%%%%%%%%%%%%%%%%%%

A more general condition is the Lindeberg condition.

For every

\[
\varepsilon>0,
\]

require

\[
\boxed{
\frac1{s_n^2}
\sum_{k=1}^{n}
\mathbb{E}
\left[
(X_k-\mu_k)^2
\mathbf{1}_{
\{
|X_k-\mu_k|
>
\varepsilon s_n
\}
}
\right]
\to0.
}
\]

This requires large individual deviations to contribute negligibly to
the total variance.

%%%%%%%%%%%%%%%%%%%%%%%%%%%%%%%%%%%%%%%%%%%%%%%%%%%%%%%%%%%%
\subsection{Lindeberg--Feller Central Limit Theorem Preview}
\label{subsec:lindeberg-feller-clt}
%%%%%%%%%%%%%%%%%%%%%%%%%%%%%%%%%%%%%%%%%%%%%%%%%%%%%%%%%%%%

The Lindeberg--Feller theorem provides a very general CLT for arrays of
independent random variables.

Under appropriate variance normalization and the Lindeberg condition,

\[
\boxed{
\frac{
\sum_k
(X_{n,k}-\mu_{n,k})
}{
s_n
}
\xrightarrow{d}
N(0,1).
}
\]

This framework handles triangular arrays.

%%%%%%%%%%%%%%%%%%%%%%%%%%%%%%%%%%%%%%%%%%%%%%%%%%%%%%%%%%%%
\subsection{Triangular Arrays}
\label{subsec:triangular-arrays}
%%%%%%%%%%%%%%%%%%%%%%%%%%%%%%%%%%%%%%%%%%%%%%%%%%%%%%%%%%%%

A triangular array has the form

\[
\begin{array}{cccccc}
X_{1,1}
\\
X_{2,1}
&
X_{2,2}
\\
X_{3,1}
&
X_{3,2}
&
X_{3,3}
\\
\vdots
&
\vdots
&
\vdots
&
\ddots
\end{array}
\]

The distribution of the variables may change with the row \(n\).

Triangular arrays are useful when studying approximations in which the
components themselves depend on sample size.

%%%%%%%%%%%%%%%%%%%%%%%%%%%%%%%%%%%%%%%%%%%%%%%%%%%%%%%%%%%%
\subsection{Lyapunov Implies Lindeberg}
\label{subsec:lyapunov-implies-lindeberg}
%%%%%%%%%%%%%%%%%%%%%%%%%%%%%%%%%%%%%%%%%%%%%%%%%%%%%%%%%%%%

Under standard assumptions, the Lyapunov condition implies the
Lindeberg condition.

Therefore Lyapunov's theorem is easier to verify in some applications
but is stronger than necessary.

This gives the hierarchy

\[
\boxed{
\text{Lyapunov condition}
\Longrightarrow
\text{Lindeberg condition}.
}
\]

%%%%%%%%%%%%%%%%%%%%%%%%%%%%%%%%%%%%%%%%%%%%%%%%%%%%%%%%%%%%
\subsection{Feller Condition Preview}
\label{subsec:feller-condition}
%%%%%%%%%%%%%%%%%%%%%%%%%%%%%%%%%%%%%%%%%%%%%%%%%%%%%%%%%%%%

A related requirement ensures that no single summand contributes a
macroscopic fraction of total variance.

A common form is

\[
\boxed{
\max_{1\leq k\leq n}
\frac{
\sigma_{n,k}^2
}{
s_n^2
}
\to0.
}
\]

This is often called an infinitesimality or Feller-type condition.

%%%%%%%%%%%%%%%%%%%%%%%%%%%%%%%%%%%%%%%%%%%%%%%%%%%%%%%%%%%%
\subsection{Berry--Esseen Theorem}
\label{subsec:berry-esseen}
%%%%%%%%%%%%%%%%%%%%%%%%%%%%%%%%%%%%%%%%%%%%%%%%%%%%%%%%%%%%

The classical CLT states convergence but does not specify the rate.

The Berry--Esseen theorem gives a quantitative bound.

Suppose \(X_i\) are i.i.d. with

\[
\mathbb{E}[X_i]=\mu,
\]

\[
\operatorname{Var}(X_i)=\sigma^2>0,
\]

and finite third absolute central moment

\[
\rho
=
\mathbb{E}[|X_i-\mu|^3].
\]

Then there exists a universal constant \(C\) such that

\[
\boxed{
\sup_x
\left|
P
\left(
\frac{
S_n-n\mu
}{
\sigma\sqrt{n}
}
\leq x
\right)
-
\Phi(x)
\right|
\leq
\frac{
C\rho
}{
\sigma^3\sqrt{n}
}.
}
\]

%%%%%%%%%%%%%%%%%%%%%%%%%%%%%%%%%%%%%%%%%%%%%%%%%%%%%%%%%%%%
\subsection{Meaning of Berry--Esseen}
\label{subsec:meaning-berry-esseen}
%%%%%%%%%%%%%%%%%%%%%%%%%%%%%%%%%%%%%%%%%%%%%%%%%%%%%%%%%%%%

The theorem shows that the approximation error is at most on the order

\[
\boxed{
\frac1{\sqrt{n}}.
}
\]

It also reveals why distributions with large third absolute moments may
converge more slowly.

Thus skewness and heavy tails can reduce finite-sample accuracy.

%%%%%%%%%%%%%%%%%%%%%%%%%%%%%%%%%%%%%%%%%%%%%%%%%%%%%%%%%%%%
\subsection{Standard Error}
\label{subsec:standard-error-clt}
%%%%%%%%%%%%%%%%%%%%%%%%%%%%%%%%%%%%%%%%%%%%%%%%%%%%%%%%%%%%

For i.i.d. observations,

\[
\operatorname{SD}(\overline{X}_n)
=
\frac{\sigma}{\sqrt{n}}.
\]

This is called the standard error of the sample mean.

Thus,

\[
\boxed{
\operatorname{SE}(\overline{X})
=
\frac{\sigma}{\sqrt{n}}.
}
\]

If \(\sigma\) is unknown, statistics often replaces it by a sample
estimate.

%%%%%%%%%%%%%%%%%%%%%%%%%%%%%%%%%%%%%%%%%%%%%%%%%%%%%%%%%%%%
\subsection{Effect of Sample Size on Standard Error}
\label{subsec:sample-size-standard-error}
%%%%%%%%%%%%%%%%%%%%%%%%%%%%%%%%%%%%%%%%%%%%%%%%%%%%%%%%%%%%

Because

\[
\operatorname{SE}
=
\frac{\sigma}{\sqrt{n}},
\]

doubling \(n\) does not halve the standard error.

To halve the standard error, one must multiply sample size by \(4\).

Indeed,

\[
\frac{
\sigma
}{
\sqrt{4n}
}
=
\frac12
\frac{
\sigma
}{
\sqrt{n}
}.
\]

Thus,

\[
\boxed{
\text{error reduction by factor }c
\text{ typically requires sample-size increase by factor }c^2.
}
\]

%%%%%%%%%%%%%%%%%%%%%%%%%%%%%%%%%%%%%%%%%%%%%%%%%%%%%%%%%%%%
\subsection{Worked Example: Required Sample Size Scaling}
\label{subsec:worked-clt-sample-size-scaling}
%%%%%%%%%%%%%%%%%%%%%%%%%%%%%%%%%%%%%%%%%%%%%%%%%%%%%%%%%%%%

\begin{workedexamplebox}{Reducing Standard Error}

Suppose a sample mean based on

\[
n=100
\]

has standard error

\[
2.
\]

To reduce the standard error to

\[
1,
\]

the error must be halved.

Therefore sample size should be multiplied by

\[
2^2=4.
\]

\begin{answerbox}
\[
\boxed{
n_{\text{new}}
=
400.
}
\]
\end{answerbox}

\end{workedexamplebox}

%%%%%%%%%%%%%%%%%%%%%%%%%%%%%%%%%%%%%%%%%%%%%%%%%%%%%%%%%%%%
\subsection{CLT and Confidence Intervals Preview}
\label{subsec:clt-confidence-interval-preview}
%%%%%%%%%%%%%%%%%%%%%%%%%%%%%%%%%%%%%%%%%%%%%%%%%%%%%%%%%%%%

The CLT implies that approximately

\[
\frac{
\overline{X}_n-\mu
}{
\sigma/\sqrt{n}
}
\sim
N(0,1)
\]

for large \(n\).

Thus,

\[
P
\left(
-z_{\alpha/2}
\leq
\frac{
\overline{X}_n-\mu
}{
\sigma/\sqrt{n}
}
\leq
z_{\alpha/2}
\right)
\]

is approximately

\[
1-\alpha.
\]

Rearranging leads to confidence intervals for \(\mu\).

This is developed systematically in Chapter 8.

%%%%%%%%%%%%%%%%%%%%%%%%%%%%%%%%%%%%%%%%%%%%%%%%%%%%%%%%%%%%
\subsection{CLT and Hypothesis Testing Preview}
\label{subsec:clt-hypothesis-testing-preview}
%%%%%%%%%%%%%%%%%%%%%%%%%%%%%%%%%%%%%%%%%%%%%%%%%%%%%%%%%%%%

Suppose a null hypothesis specifies

\[
\mu=\mu_0.
\]

Then a standardized sample mean

\[
\boxed{
Z
=
\frac{
\overline{X}_n-\mu_0
}{
\sigma/\sqrt{n}
}
}
\]

is approximately standard normal for large samples under suitable
conditions.

This provides the basis for large-sample \(z\)-tests.

%%%%%%%%%%%%%%%%%%%%%%%%%%%%%%%%%%%%%%%%%%%%%%%%%%%%%%%%%%%%
\subsection{CLT and Monte Carlo Error}
\label{subsec:clt-monte-carlo-error}
%%%%%%%%%%%%%%%%%%%%%%%%%%%%%%%%%%%%%%%%%%%%%%%%%%%%%%%%%%%%

Suppose

\[
Y_1,\ldots,Y_n
\]

are i.i.d. simulation outputs with

\[
\mathbb{E}[Y_i]=\mu
\]

and finite variance

\[
\sigma^2.
\]

The Monte Carlo estimator is

\[
\widehat{\mu}_n
=
\frac1n
\sum_{i=1}^{n}Y_i.
\]

The LLN says

\[
\widehat{\mu}_n
\to
\mu.
\]

The CLT additionally gives

\[
\boxed{
\frac{
\widehat{\mu}_n-\mu
}{
\sigma/\sqrt{n}
}
\approx
N(0,1).
}
\]

Thus the CLT quantifies simulation error.

%%%%%%%%%%%%%%%%%%%%%%%%%%%%%%%%%%%%%%%%%%%%%%%%%%%%%%%%%%%%
\subsection{Worked Example: Monte Carlo Error Probability}
\label{subsec:worked-monte-carlo-clt}
%%%%%%%%%%%%%%%%%%%%%%%%%%%%%%%%%%%%%%%%%%%%%%%%%%%%%%%%%%%%

\begin{workedexamplebox}{Approximate Monte Carlo Error}

Suppose a Monte Carlo output has

\[
\sigma=5.
\]

Using

\[
n=10{,}000
\]

independent simulations,

\[
\operatorname{SE}
=
\frac5{100}
=
0.05.
\]

Therefore,

\[
\widehat{\mu}
\]

is approximately normal around \(\mu\) with standard deviation \(0.05\).

\begin{answerbox}
\[
\boxed{
\widehat{\mu}
\approx
N(\mu,0.05^2).
}
\]
\end{answerbox}

\end{workedexamplebox}

%%%%%%%%%%%%%%%%%%%%%%%%%%%%%%%%%%%%%%%%%%%%%%%%%%%%%%%%%%%%
\subsection{CLT and Measurement Error}
\label{subsec:clt-measurement-error}
%%%%%%%%%%%%%%%%%%%%%%%%%%%%%%%%%%%%%%%%%%%%%%%%%%%%%%%%%%%%

Suppose an observed quantity is the sum of many small approximately
independent effects:

\[
E
=
E_1+\cdots+E_n.
\]

Even when the individual components are not normal, the combined error
may be approximately normal after suitable centering and scaling.

This helps explain the frequent appearance of Gaussian measurement
models.

%%%%%%%%%%%%%%%%%%%%%%%%%%%%%%%%%%%%%%%%%%%%%%%%%%%%%%%%%%%%
\subsection{Universality of the Normal Limit}
\label{subsec:normal-limit-universality}
%%%%%%%%%%%%%%%%%%%%%%%%%%%%%%%%%%%%%%%%%%%%%%%%%%%%%%%%%%%%

One of the remarkable features of the CLT is that many different parent
distributions lead to the same limiting standard normal law.

For example, suitable sums of:

\[
\text{Bernoulli variables},
\]

\[
\text{uniform variables},
\]

\[
\text{exponential variables},
\]

and many other finite-variance distributions all converge after
standardization to

\[
N(0,1).
\]

Thus the normal law is a universal limit under broad conditions.

%%%%%%%%%%%%%%%%%%%%%%%%%%%%%%%%%%%%%%%%%%%%%%%%%%%%%%%%%%%%
\subsection{What the Limit Forgets}
\label{subsec:clt-forgets-parent-distribution}
%%%%%%%%%%%%%%%%%%%%%%%%%%%%%%%%%%%%%%%%%%%%%%%%%%%%%%%%%%%%

The limiting standard normal distribution depends on the original
distribution primarily through:

\[
\boxed{
\mu
}
\]

and

\[
\boxed{
\sigma^2.
}
\]

After centering and scaling, much of the detailed shape of the original
distribution disappears asymptotically.

This is one reason normal approximations are so widely applicable.

%%%%%%%%%%%%%%%%%%%%%%%%%%%%%%%%%%%%%%%%%%%%%%%%%%%%%%%%%%%%
\subsection{Skewness Under Averaging}
\label{subsec:skewness-under-averaging}
%%%%%%%%%%%%%%%%%%%%%%%%%%%%%%%%%%%%%%%%%%%%%%%%%%%%%%%%%%%%

Suppose i.i.d. variables have finite third cumulant

\[
\kappa_3.
\]

For the sum \(S_n\),

\[
\kappa_3(S_n)
=
n\kappa_3.
\]

The variance is

\[
n\sigma^2.
\]

Therefore standardized skewness scales like

\[
\boxed{
\frac1{\sqrt{n}}.
}
\]

Thus skewness tends to disappear as independent observations are
averaged.

%%%%%%%%%%%%%%%%%%%%%%%%%%%%%%%%%%%%%%%%%%%%%%%%%%%%%%%%%%%%
\subsection{Higher Cumulants and Normality}
\label{subsec:higher-cumulants-clt}
%%%%%%%%%%%%%%%%%%%%%%%%%%%%%%%%%%%%%%%%%%%%%%%%%%%%%%%%%%%%

For independent sums, cumulants add.

After normalization by \(\sqrt{n}\), higher standardized cumulants tend
to shrink.

For order \(r\geq3\), the standardized cumulant typically scales as

\[
\boxed{
n^{1-r/2}.
}
\]

Since

\[
1-\frac r2<0
\]

for \(r>2\), these terms vanish.

The first two cumulants remain as mean and variance.

This gives an important heuristic explanation for the normal limit.

%%%%%%%%%%%%%%%%%%%%%%%%%%%%%%%%%%%%%%%%%%%%%%%%%%%%%%%%%%%%
\subsection{Local Central Limit Theorem Preview}
\label{subsec:local-clt-preview}
%%%%%%%%%%%%%%%%%%%%%%%%%%%%%%%%%%%%%%%%%%%%%%%%%%%%%%%%%%%%

The ordinary CLT approximates distribution functions.

A local central limit theorem gives approximations to individual
probabilities or densities.

For suitable lattice-valued sums, one may obtain approximations of the
form

\[
\boxed{
P(S_n=k)
\approx
\frac{
1
}{
\sigma\sqrt{2\pi n}
}
\exp
\left(
-
\frac{
(k-n\mu)^2
}{
2n\sigma^2
}
\right).
}
\]

This is stronger than ordinary distributional convergence.

%%%%%%%%%%%%%%%%%%%%%%%%%%%%%%%%%%%%%%%%%%%%%%%%%%%%%%%%%%%%
\subsection{Edgeworth Expansions Preview}
\label{subsec:edgeworth-expansions-preview}
%%%%%%%%%%%%%%%%%%%%%%%%%%%%%%%%%%%%%%%%%%%%%%%%%%%%%%%%%%%%

The normal approximation can be refined using higher moments or
cumulants.

An Edgeworth expansion adds correction terms involving quantities such
as skewness and kurtosis.

Symbolically,

\[
\boxed{
F_n(x)
=
\Phi(x)
+
\text{correction terms}
+
\text{smaller error}.
}
\]

These expansions improve finite-sample approximations under suitable
regularity conditions.

%%%%%%%%%%%%%%%%%%%%%%%%%%%%%%%%%%%%%%%%%%%%%%%%%%%%%%%%%%%%
\subsection{Stable Laws Preview}
\label{subsec:stable-laws-clt-preview}
%%%%%%%%%%%%%%%%%%%%%%%%%%%%%%%%%%%%%%%%%%%%%%%%%%%%%%%%%%%%

The normal distribution is a stable distribution.

A distribution is stable if sums of independent copies, after suitable
location and scale changes, remain in the same distribution family.

When variance is infinite, nonnormal stable distributions may replace
the Gaussian as the limiting law.

Thus the classical CLT is part of a larger theory of stable limits.

%%%%%%%%%%%%%%%%%%%%%%%%%%%%%%%%%%%%%%%%%%%%%%%%%%%%%%%%%%%%
\subsection{Generalized Central Limit Theorem Preview}
\label{subsec:generalized-clt-preview}
%%%%%%%%%%%%%%%%%%%%%%%%%%%%%%%%%%%%%%%%%%%%%%%%%%%%%%%%%%%%

The generalized central limit theorem characterizes possible
nondegenerate limits of properly normalized sums of i.i.d. variables.

The possible limits are stable distributions.

The normal distribution corresponds to the finite-variance case.

Heavy-tailed models can instead lead to stable laws with power-law
tails.

%%%%%%%%%%%%%%%%%%%%%%%%%%%%%%%%%%%%%%%%%%%%%%%%%%%%%%%%%%%%
\subsection{Multivariate Central Limit Theorem}
\label{subsec:multivariate-clt}
%%%%%%%%%%%%%%%%%%%%%%%%%%%%%%%%%%%%%%%%%%%%%%%%%%%%%%%%%%%%

Suppose

\[
\mathbf{X}_1,\mathbf{X}_2,\ldots
\]

are i.i.d. random vectors in

\[
\mathbb{R}^d
\]

with mean vector

\[
\boldsymbol{\mu}
\]

and finite covariance matrix

\[
\Sigma.
\]

Define

\[
\overline{\mathbf{X}}_n
=
\frac1n
\sum_{i=1}^{n}\mathbf{X}_i.
\]

Then under the multivariate CLT,

\[
\boxed{
\sqrt{n}
\left(
\overline{\mathbf{X}}_n
-
\boldsymbol{\mu}
\right)
\xrightarrow{d}
N_d(\mathbf{0},\Sigma).
}
\]

%%%%%%%%%%%%%%%%%%%%%%%%%%%%%%%%%%%%%%%%%%%%%%%%%%%%%%%%%%%%
\subsection{Meaning of the Multivariate CLT}
\label{subsec:meaning-multivariate-clt}
%%%%%%%%%%%%%%%%%%%%%%%%%%%%%%%%%%%%%%%%%%%%%%%%%%%%%%%%%%%%

The scalar CLT describes one average.

The multivariate CLT describes several averages simultaneously.

The covariance matrix

\[
\Sigma
\]

captures dependence among the coordinates.

Thus the limiting shape is a multivariate normal distribution rather
than several unrelated scalar normals.

%%%%%%%%%%%%%%%%%%%%%%%%%%%%%%%%%%%%%%%%%%%%%%%%%%%%%%%%%%%%
\subsection{Cramér--Wold Device Preview}
\label{subsec:cramer-wold-preview}
%%%%%%%%%%%%%%%%%%%%%%%%%%%%%%%%%%%%%%%%%%%%%%%%%%%%%%%%%%%%

A useful principle states that convergence of random vectors can be
checked through all linear combinations.

For a vector

\[
\mathbf{a},
\]

consider

\[
\mathbf{a}^T
\mathbf{X}_n.
\]

The Cramér--Wold device says, roughly, that convergence of all such
linear combinations can characterize multivariate convergence in
distribution.

This reduces multivariate CLT proofs to scalar CLTs.

%%%%%%%%%%%%%%%%%%%%%%%%%%%%%%%%%%%%%%%%%%%%%%%%%%%%%%%%%%%%
\subsection{Delta Method Preview}
\label{subsec:delta-method-preview}
%%%%%%%%%%%%%%%%%%%%%%%%%%%%%%%%%%%%%%%%%%%%%%%%%%%%%%%%%%%%

Suppose

\[
\sqrt{n}
(T_n-\theta)
\xrightarrow{d}
N(0,\sigma^2)
\]

and \(g\) is differentiable at \(\theta\).

Then, under suitable conditions,

\[
\boxed{
\sqrt{n}
\left(
g(T_n)-g(\theta)
\right)
\xrightarrow{d}
N
\left(
0,
[g'(\theta)]^2\sigma^2
\right).
}
\]

This is called the delta method.

It transfers asymptotic normality through smooth transformations.

%%%%%%%%%%%%%%%%%%%%%%%%%%%%%%%%%%%%%%%%%%%%%%%%%%%%%%%%%%%%
\subsection{Worked Example: Delta Method}
\label{subsec:worked-delta-method}
%%%%%%%%%%%%%%%%%%%%%%%%%%%%%%%%%%%%%%%%%%%%%%%%%%%%%%%%%%%%

\begin{workedexamplebox}{Transforming an Asymptotically Normal Estimator}

Suppose

\[
\sqrt{n}
(T_n-\theta)
\xrightarrow{d}
N(0,\sigma^2)
\]

and

\[
g(x)=x^2.
\]

Then

\[
g'(\theta)=2\theta.
\]

Therefore,

\begin{answerbox}
\[
\boxed{
\sqrt{n}
(T_n^2-\theta^2)
\xrightarrow{d}
N
\left(
0,
4\theta^2\sigma^2
\right).
}
\]
\end{answerbox}

\end{workedexamplebox}

%%%%%%%%%%%%%%%%%%%%%%%%%%%%%%%%%%%%%%%%%%%%%%%%%%%%%%%%%%%%
\subsection{Slutsky's Theorem}
\label{subsec:slutsky-clt}
%%%%%%%%%%%%%%%%%%%%%%%%%%%%%%%%%%%%%%%%%%%%%%%%%%%%%%%%%%%%

A fundamental asymptotic result is Slutsky's theorem.

Suppose

\[
X_n
\xrightarrow{d}
X
\]

and

\[
Y_n
\xrightarrow{P}
c
\]

for a constant \(c\).

Then

\[
\boxed{
X_n+Y_n
\xrightarrow{d}
X+c,
}
\]

\[
\boxed{
X_nY_n
\xrightarrow{d}
cX,
}
\]

and if

\[
c\neq0,
\]

\[
\boxed{
\frac{X_n}{Y_n}
\xrightarrow{d}
\frac{X}{c}.
}
\]

%%%%%%%%%%%%%%%%%%%%%%%%%%%%%%%%%%%%%%%%%%%%%%%%%%%%%%%%%%%%
\subsection{Replacing Unknown Standard Deviations}
\label{subsec:replace-sigma-slutsky}
%%%%%%%%%%%%%%%%%%%%%%%%%%%%%%%%%%%%%%%%%%%%%%%%%%%%%%%%%%%%

Suppose the CLT gives

\[
\frac{
\sqrt{n}
(\overline{X}_n-\mu)
}{
\sigma
}
\xrightarrow{d}
N(0,1).
\]

If an estimator \(S_n\) satisfies

\[
S_n
\xrightarrow{P}
\sigma,
\]

then Slutsky's theorem gives

\[
\boxed{
\frac{
\sqrt{n}
(\overline{X}_n-\mu)
}{
S_n
}
\xrightarrow{d}
N(0,1).
}
\]

This is fundamental for large-sample statistical inference.

%%%%%%%%%%%%%%%%%%%%%%%%%%%%%%%%%%%%%%%%%%%%%%%%%%%%%%%%%%%%
\subsection{Studentization Preview}
\label{subsec:studentization-preview}
%%%%%%%%%%%%%%%%%%%%%%%%%%%%%%%%%%%%%%%%%%%%%%%%%%%%%%%%%%%%

Replacing an unknown standard deviation by an estimated one is called
studentization.

A statistic of the form

\[
\boxed{
\frac{
\text{estimate}-\text{parameter}
}{
\text{estimated standard error}
}
}
\]

often has an asymptotically standard normal distribution.

For normal finite samples, special exact \(t\)-distribution results are
available.

%%%%%%%%%%%%%%%%%%%%%%%%%%%%%%%%%%%%%%%%%%%%%%%%%%%%%%%%%%%%
\subsection{CLT for Dependent Variables Preview}
\label{subsec:dependent-clt-preview}
%%%%%%%%%%%%%%%%%%%%%%%%%%%%%%%%%%%%%%%%%%%%%%%%%%%%%%%%%%%%

Independence can be weakened in many central limit theorems.

CLTs exist for:

\[
\text{mixing sequences},
\]

\[
\text{Markov chains},
\]

\[
\text{martingales},
\]

and

\[
\text{stationary stochastic processes}.
\]

Dependence changes the effective asymptotic variance.

%%%%%%%%%%%%%%%%%%%%%%%%%%%%%%%%%%%%%%%%%%%%%%%%%%%%%%%%%%%%
\subsection{Long-Run Variance Preview}
\label{subsec:long-run-variance-preview}
%%%%%%%%%%%%%%%%%%%%%%%%%%%%%%%%%%%%%%%%%%%%%%%%%%%%%%%%%%%%

For a stationary dependent sequence with covariance function

\[
\gamma(k)
=
\operatorname{Cov}(X_0,X_k),
\]

the asymptotic variance of a sample mean may involve

\[
\boxed{
\gamma(0)
+
2
\sum_{k=1}^{\infty}
\gamma(k).
}
\]

Thus serial dependence can increase or decrease long-run variance.

This becomes important in time-series analysis.

%%%%%%%%%%%%%%%%%%%%%%%%%%%%%%%%%%%%%%%%%%%%%%%%%%%%%%%%%%%%
\subsection{Martingale Central Limit Theorem Preview}
\label{subsec:martingale-clt-preview}
%%%%%%%%%%%%%%%%%%%%%%%%%%%%%%%%%%%%%%%%%%%%%%%%%%%%%%%%%%%%

Martingale differences provide another important dependent framework.

Under suitable conditional variance and regularity conditions, sums of
martingale differences can converge to normal distributions.

These results are widely used in:

\[
\text{stochastic processes},
\]

\[
\text{econometrics},
\]

and

\[
\text{modern asymptotic statistics}.
\]

%%%%%%%%%%%%%%%%%%%%%%%%%%%%%%%%%%%%%%%%%%%%%%%%%%%%%%%%%%%%
\subsection{CLT Workflow}
\label{subsec:clt-workflow}
%%%%%%%%%%%%%%%%%%%%%%%%%%%%%%%%%%%%%%%%%%%%%%%%%%%%%%%%%%%%

When using a central limit theorem:

\begin{enumerate}
    \item identify the random sum or average;
    \item determine the mean;
    \item determine the variance;
    \item check independence or other applicable assumptions;
    \item check whether a finite-variance CLT is appropriate;
    \item center by subtracting the mean;
    \item scale by the correct standard deviation;
    \item convert the target event into a standardized normal event;
    \item use \(\Phi\) to evaluate the approximation;
    \item apply a continuity correction when approximating discrete
          variables;
    \item interpret the approximation rather than treating it as exact.
\end{enumerate}

%%%%%%%%%%%%%%%%%%%%%%%%%%%%%%%%%%%%%%%%%%%%%%%%%%%%%%%%%%%%
\subsection{Common Errors}
\label{subsec:clt-common-errors}
%%%%%%%%%%%%%%%%%%%%%%%%%%%%%%%%%%%%%%%%%%%%%%%%%%%%%%%%%%%%

\subsubsection{Confusing the CLT with the LLN}

The LLN says the average converges to its mean.

The CLT describes the distribution of the scaled error around that mean.

%%%%%%%%%%%%%%%%%%%%%%%%%%%%%%%%%%%%%%%%%%%%%%%%%%%%%%%%%%%%
\subsubsection{Forgetting to Center}

The CLT standardization requires subtracting

\[
n\mu
\]

for a sum or

\[
\mu
\]

for a sample mean.

%%%%%%%%%%%%%%%%%%%%%%%%%%%%%%%%%%%%%%%%%%%%%%%%%%%%%%%%%%%%
\subsubsection{Using the Wrong Scaling}

For a sum,

\[
\operatorname{SD}(S_n)
=
\sigma\sqrt{n}.
\]

For a sample mean,

\[
\operatorname{SD}(\overline{X}_n)
=
\frac{\sigma}{\sqrt{n}}.
\]

These must not be confused.

%%%%%%%%%%%%%%%%%%%%%%%%%%%%%%%%%%%%%%%%%%%%%%%%%%%%%%%%%%%%
\subsubsection{Dividing by Variance Instead of Standard Deviation}

Standardization uses

\[
\sigma
\]

or

\[
\sigma/\sqrt{n},
\]

not

\[
\sigma^2.
\]

%%%%%%%%%%%%%%%%%%%%%%%%%%%%%%%%%%%%%%%%%%%%%%%%%%%%%%%%%%%%
\subsubsection{Assuming the Sample Must Be Normal}

The classical CLT does not require the original population to be
normal.

It requires appropriate moment and independence assumptions.

%%%%%%%%%%%%%%%%%%%%%%%%%%%%%%%%%%%%%%%%%%%%%%%%%%%%%%%%%%%%
\subsubsection{Assuming the Approximation Is Exact}

Unless the underlying distribution gives exact normal sums, the CLT is
an approximation for finite \(n\).

%%%%%%%%%%%%%%%%%%%%%%%%%%%%%%%%%%%%%%%%%%%%%%%%%%%%%%%%%%%%
\subsubsection{Using a Fixed Rule Such as \(n=30\) Universally}

No single sample-size rule works for every distribution.

Highly skewed or heavy-tailed populations may require much larger
samples.

%%%%%%%%%%%%%%%%%%%%%%%%%%%%%%%%%%%%%%%%%%%%%%%%%%%%%%%%%%%%
\subsubsection{Ignoring Infinite Variance}

The classical finite-variance CLT may fail for heavy-tailed
distributions with infinite variance.

%%%%%%%%%%%%%%%%%%%%%%%%%%%%%%%%%%%%%%%%%%%%%%%%%%%%%%%%%%%%
\subsubsection{Forgetting Continuity Correction}

When a discrete distribution is approximated by a continuous normal
distribution, a \(0.5\) correction often improves the approximation.

%%%%%%%%%%%%%%%%%%%%%%%%%%%%%%%%%%%%%%%%%%%%%%%%%%%%%%%%%%%%
\subsubsection{Applying Continuity Correction to Already Continuous Data}

Continuity correction is intended for discrete-to-continuous
approximations.

It is not needed for ordinary continuous sample means.

%%%%%%%%%%%%%%%%%%%%%%%%%%%%%%%%%%%%%%%%%%%%%%%%%%%%%%%%%%%%
\subsubsection{Confusing Standard Error with Population Standard Deviation}

The population standard deviation is

\[
\sigma.
\]

The standard error of an independent sample mean is

\[
\boxed{
\frac{\sigma}{\sqrt{n}}.
}
\]

%%%%%%%%%%%%%%%%%%%%%%%%%%%%%%%%%%%%%%%%%%%%%%%%%%%%%%%%%%%%
\subsubsection{Assuming Independence When It Is Not Present}

Dependent data may require a different CLT and a different asymptotic
variance.

%%%%%%%%%%%%%%%%%%%%%%%%%%%%%%%%%%%%%%%%%%%%%%%%%%%%%%%%%%%%
\subsection{Applications}
\label{subsec:clt-applications}
%%%%%%%%%%%%%%%%%%%%%%%%%%%%%%%%%%%%%%%%%%%%%%%%%%%%%%%%%%%%

\begin{applicationbox}

\textbf{Statistics.}
The CLT provides the basis for large-sample confidence intervals,
hypothesis tests, and asymptotic distributions of estimators.

\medskip

\textbf{Survey Sampling.}
Sample proportions and means are often approximated by normal
distributions.

\medskip

\textbf{Quality Control.}
Counts, defect proportions, and average measurements can be analyzed
using normal approximations.

\medskip

\textbf{Monte Carlo Methods.}
The CLT quantifies the distribution of simulation error around an
estimated expectation.

\medskip

\textbf{Finance.}
Aggregated returns and portfolio quantities may exhibit approximate
normal behavior under suitable assumptions, though dependence and heavy
tails must be considered carefully.

\medskip

\textbf{Engineering.}
Measurements formed from many small independent effects often exhibit
approximately Gaussian errors.

\medskip

\textbf{Operations Research.}
Aggregate demand, service quantities, and workloads may be approximated
using CLT methods.

\medskip

\textbf{Stochastic Processes.}
CLTs describe fluctuations of random walks, counting processes, Markov
chains, and time averages around their long-run limits.

\end{applicationbox}

%%%%%%%%%%%%%%%%%%%%%%%%%%%%%%%%%%%%%%%%%%%%%%%%%%%%%%%%%%%%
\subsection{Exercises}
\label{subsec:clt-exercises}
%%%%%%%%%%%%%%%%%%%%%%%%%%%%%%%%%%%%%%%%%%%%%%%%%%%%%%%%%%%%

\begin{exercise}[1]
State the classical i.i.d. central limit theorem.
\end{exercise}

\begin{exercise}[1]
Define convergence in distribution.
\end{exercise}

\begin{exercise}[1]
State the standardized form of a sample sum.
\end{exercise}

\begin{exercise}[1]
State the standardized form of a sample mean.
\end{exercise}

\begin{exercise}[1]
Explain the difference between the LLN and CLT.
\end{exercise}

\begin{exercise}[1]
What is the standard error of a sample mean based on \(n\) independent
observations with population standard deviation \(\sigma\)?
\end{exercise}

\begin{exercise}[1]
State the de Moivre--Laplace theorem.
\end{exercise}

\begin{exercise}[1]
Explain why continuity correction is used.
\end{exercise}

\begin{exercise}[2]
Suppose

\[
\mu=40,
\qquad
\sigma=10,
\qquad
n=25.
\]

Find the approximate distribution of

\[
\overline{X}_n.
\]
\end{exercise}

\begin{exercise}[2]
Using the previous information, approximate

\[
P(\overline{X}_n\leq44).
\]
\end{exercise}

\begin{exercise}[2]
Suppose

\[
\mu=100,
\qquad
\sigma=20,
\qquad
n=100.
\]

Find the standard error of the sample mean.
\end{exercise}

\begin{exercise}[2]
For the previous problem, approximately find

\[
P(98<\overline{X}_n<102).
\]
\end{exercise}

\begin{exercise}[2]
Let

\[
X\sim\operatorname{Bin}(200,0.4).
\]

Find the mean and variance of \(X\).
\end{exercise}

\begin{exercise}[2]
Use a normal approximation with continuity correction to approximate

\[
P(X\leq90)
\]

for the previous binomial variable.
\end{exercise}

\begin{exercise}[2]
Use a normal approximation to estimate

\[
P(X=80)
\]

for the same binomial variable.
\end{exercise}

\begin{exercise}[2]
Let

\[
X\sim\operatorname{Poisson}(100).
\]

Approximate

\[
P(X\leq110)
\]

using a normal approximation and continuity correction.
\end{exercise}

\begin{exercise}[2]
Suppose

\[
\widehat{p}
\]

is based on

\[
n=500
\]

Bernoulli trials with

\[
p=0.2.
\]

Find the approximate mean and standard deviation of \(\widehat{p}\).
\end{exercise}

\begin{exercise}[2]
For the previous proportion, approximate

\[
P(\widehat{p}>0.23).
\]
\end{exercise}

\begin{exercise}[3]
Show that

\[
\frac{
S_n-n\mu
}{
\sigma\sqrt{n}
}
=
\frac{
\sqrt{n}
(\overline{X}_n-\mu)
}{
\sigma
}.
\]
\end{exercise}

\begin{exercise}[3]
Prove that the standardized sum has mean \(0\) and variance \(1\).
\end{exercise}

\begin{exercise}[3]
Starting with the characteristic-function expansion

\[
\varphi_X(t)
=
1-\frac{t^2}{2}+o(t^2),
\]

derive the main limiting expression in the characteristic-function
proof of the CLT.
\end{exercise}

\begin{exercise}[3]
Show that

\[
\left(
1-\frac{t^2}{2n}
\right)^n
\to
e^{-t^2/2}.
\]
\end{exercise}

\begin{exercise}[3]
Use characteristic functions to outline a proof of the classical CLT.
\end{exercise}

\begin{exercise}[3]
Derive the de Moivre--Laplace theorem from the classical CLT by viewing
a binomial variable as a sum of Bernoulli variables.
\end{exercise}

\begin{exercise}[3]
Explain mathematically why the standard deviation of a sample mean
shrinks at rate

\[
n^{-1/2}.
\]
\end{exercise}

\begin{exercise}[3]
How much must sample size increase to reduce standard error by a factor
of \(3\)?
\end{exercise}

\begin{exercise}[3]
Suppose a population has finite third absolute central moment. Explain
what the Berry--Esseen theorem says about normal-approximation error.
\end{exercise}

\begin{exercise}[3]
State the Lyapunov condition.
\end{exercise}

\begin{exercise}[3]
Explain conceptually why the Lyapunov condition prevents one variable
from dominating the total sum.
\end{exercise}

\begin{exercise}[3]
State the Lindeberg condition.
\end{exercise}

\begin{exercise}[3]
Compare the Lindeberg--Lévy, Lyapunov, and Lindeberg--Feller central
limit theorems.
\end{exercise}

\begin{exercise}[3]
Suppose

\[
X_1,\ldots,X_n
\]

are independent exponential variables with finite variance.

Explain why their standardized average satisfies the classical CLT when
they are identically distributed.
\end{exercise}

\begin{exercise}[3]
Explain why the standard Cauchy distribution does not satisfy the
classical finite-variance CLT.
\end{exercise}

\begin{exercise}[3]
Suppose

\[
X_i
\]

are i.i.d. with finite variance and

\[
Y_i=aX_i+b.
\]

Show that the standardized sample mean of the \(Y_i\) has the same CLT
limit after using the appropriate mean and variance.
\end{exercise}

\begin{exercise}[3]
Use the CLT to approximate the probability that a sum of \(100\)
independent uniform \((0,1)\) variables exceeds \(55\).
\end{exercise}

\begin{exercise}[3]
Use the CLT to approximate the distribution of the average of \(64\)
independent exponential variables with rate \(2\).
\end{exercise}

\begin{exercise}[3]
Explain why a sample mean from a normal population is exactly normal
rather than only asymptotically normal.
\end{exercise}

\begin{exercise}[4]
Study the Berry--Esseen theorem and current known bounds on the universal
constant.
\end{exercise}

\begin{exercise}[4]
Prove the Lyapunov central limit theorem using characteristic functions.
\end{exercise}

\begin{exercise}[4]
Study the Lindeberg--Feller theorem for triangular arrays.
\end{exercise}

\begin{exercise}[4]
Investigate the Feller condition and its relationship with the
Lindeberg condition.
\end{exercise}

\begin{exercise}[4]
Study a local central limit theorem for lattice random variables.
\end{exercise}

\begin{exercise}[4]
Investigate Edgeworth expansions and derive a first skewness correction
to the normal approximation.
\end{exercise}

\begin{exercise}[4]
Study stable distributions and the generalized central limit theorem.
\end{exercise}

\begin{exercise}[4]
Investigate multivariate central limit theorems.
\end{exercise}

\begin{exercise}[4]
Study the Cramér--Wold device and use it to prove a multivariate CLT.
\end{exercise}

\begin{exercise}[4]
Study the delta method and derive asymptotic distributions for smooth
functions of estimators.
\end{exercise}

\begin{exercise}[4]
Investigate Slutsky's theorem and its use in studentized statistics.
\end{exercise}

\begin{exercise}[4]
Study central limit theorems for martingale difference sequences.
\end{exercise}

\begin{exercise}[4]
Investigate central limit theorems for stationary dependent sequences.
\end{exercise}

\begin{exercise}[4]
Study long-run variance estimation in time-series central limit
theorems.
\end{exercise}

%%%%%%%%%%%%%%%%%%%%%%%%%%%%%%%%%%%%%%%%%%%%%%%%%%%%%%%%%%%%
\subsection{Conceptual Summary}
\label{subsec:clt-summary}
%%%%%%%%%%%%%%%%%%%%%%%%%%%%%%%%%%%%%%%%%%%%%%%%%%%%%%%%%%%%

For i.i.d. variables with

\[
\mathbb{E}[X_i]=\mu
\]

and

\[
0<
\operatorname{Var}(X_i)
=
\sigma^2
<
\infty,
\]

the classical central limit theorem states

\[
\boxed{
\frac{
S_n-n\mu
}{
\sigma\sqrt{n}
}
\xrightarrow{d}
N(0,1).
}
\]

Equivalently,

\[
\boxed{
\frac{
\overline{X}_n-\mu
}{
\sigma/\sqrt{n}
}
\xrightarrow{d}
N(0,1).
}
\]

Therefore,

\[
\boxed{
S_n
\approx
N(n\mu,n\sigma^2)
}
\]

and

\[
\boxed{
\overline{X}_n
\approx
N
\left(
\mu,
\frac{\sigma^2}{n}
\right)
}
\]

for sufficiently large \(n\).

For a binomial random variable,

\[
\boxed{
X
\sim
\operatorname{Bin}(n,p)
\approx
N
\left(
np,
np(1-p)
\right).
}
\]

Continuity correction adjusts discrete boundaries by

\[
\boxed{
0.5.
}
\]

For sample proportions,

\[
\boxed{
\widehat{p}
\approx
N
\left(
p,
\frac{
p(1-p)
}{n}
\right).
}
\]

The LLN tells us

\[
\boxed{
\overline{X}_n-\mu
\to0,
}
\]

while the CLT explains the scale and distribution of

\[
\boxed{
\sqrt{n}
(\overline{X}_n-\mu).
}
\]

More general CLTs replace identical distribution with conditions such as

\[
\boxed{
\text{Lyapunov}
}
\]

or

\[
\boxed{
\text{Lindeberg}.
}
\]

Berry--Esseen provides a quantitative convergence rate of order

\[
\boxed{
n^{-1/2}.
}
\]

Multivariate central limit theorems extend asymptotic normality to random
vectors.

Central limit theory therefore explains why the normal distribution
appears throughout

\[
\boxed{
\text{statistics}
+
\text{simulation}
+
\text{measurement}
+
\text{stochastic modeling}.
}
\]

%%%%%%%%%%%%%%%%%%%%%%%%%%%%%%%%%%%%%%%%%%%%%%%%%%%%%%%%%%%%
\subsection{Section Connections}
\label{subsec:clt-connections}
%%%%%%%%%%%%%%%%%%%%%%%%%%%%%%%%%%%%%%%%%%%%%%%%%%%%%%%%%%%%

\chapterconnection{
Central Limit Theorems refine the Laws of Large Numbers from
Section~\ref{sec:laws-large-numbers} by describing the random
fluctuations that remain after sample averages concentrate near their
expectations. The characteristic-function machinery of
Section~\ref{sec:probability-transforms} provides a standard proof method
for the classical CLT, while the moments and covariance structures of
Section~\ref{sec:expectation-moments} determine the centering and
scaling. The next section, Measure-Theoretic Probability, develops the
formal probability-space and convergence framework underlying these
results. Chapter 8 Statistics uses central-limit theory extensively for
sampling distributions, confidence intervals, hypothesis testing,
regression, and asymptotic estimation.
}

%%%%%%%%%%%%%%%%%%%%%%%%%%%%%%%%%%%%%%%%%%%%%%%%%%%%%%%%%%%%
% SECTION SOURCE TRACKING
%%%%%%%%%%%%%%%%%%%%%%%%%%%%%%%%%%%%%%%%%%%%%%%%%%%%%%%%%%%%

%------------------------------------------------------------
% 7.10 Central Limit Theorems
%------------------------------------------------------------
%
% Source basis:
%   - Standard classical central limit theory.
%   - Standard normal approximation methods.
%   - Standard characteristic-function methods.
%   - Standard triangular-array central limit theory.
%
% Used for:
%   - convergence in distribution;
%   - standardized sums;
%   - standardized sample means;
%   - Lindeberg--Lévy central limit theorem;
%   - CDF interpretation of CLT convergence;
%   - approximate normality of sums;
%   - approximate normality of sample means;
%   - exact versus asymptotic normality;
%   - LLN versus CLT;
%   - characteristic-function proof idea;
%   - Lévy continuity theorem connection;
%   - de Moivre--Laplace theorem;
%   - normal approximation to binomial distributions;
%   - continuity corrections;
%   - Poisson normal approximation;
%   - sample proportion CLT;
%   - finite-sample approximation quality;
%   - finite-variance requirement;
%   - heavy-tailed counterexamples;
%   - Lyapunov CLT;
%   - Lyapunov condition;
%   - Lindeberg condition;
%   - Lindeberg--Feller theorem preview;
%   - triangular arrays;
%   - Feller condition preview;
%   - Berry--Esseen theorem;
%   - standard errors;
%   - CLT-based Monte Carlo error;
%   - cumulant interpretation;
%   - local CLT preview;
%   - Edgeworth expansion preview;
%   - generalized CLT and stable laws preview;
%   - multivariate CLT;
%   - Cramér--Wold device preview;
%   - delta method preview;
%   - Slutsky's theorem;
%   - studentization preview;
%   - dependent-data CLTs;
%   - long-run variance preview;
%   - martingale CLT preview;
%   - applications and exercises.
%
% Notes:
%   - Measure-Theoretic Probability is developed next in Section 7.11.
%   - Statistical applications of CLTs are developed extensively in
%     Chapter 8.
%   - Stable distributions and generalized CLTs are introduced only as
%     advanced extensions here.
%   - Dependent-data CLTs are revisited in stochastic-process and
%     time-series material.
%
%%%%%%%%%%%%%%%%%%%%%%%%%%%%%%%%%%%%%%%%%%%%%%%%%%%%%%%%%%%%